% BANKS-SOURCE: pdf=77 print=67 kind=implicit-exercise id=I-CH06-002
\begin{exercise}[A working derivation of Diracology]{I-CH06-002}
Work in four-dimensional Minkowski space with
\(\eta=\operatorname{diag}(-1,1,1,1)\), \(\epsilon^{0123}=1\), and
\(\gamma_5=\ii\gamma^0\gamma^1\gamma^2\gamma^3\).  Derive the results
collected in Appendix~C:
\begin{enumerate}[label=(\alph*)]
 \item Construct the Weyl matrices from
 \(\sigma^\mu=(\mathbf1,\boldsymbol\sigma)\) and
 \(\bar\sigma^\mu=(\mathbf1,-\boldsymbol\sigma)\).  Obtain the Dirac and
 Majorana bases by unitary changes of basis, and verify the Hermiticity
 relation \((\gamma^\mu)^\dagger=\gamma^0\gamma^\mu\gamma^0\).  Explain
 why the Majorana-basis Dirac equation separates into equations for the
 real and imaginary parts of a spinor.
 \item Show that antisymmetrized gamma products of ranks zero through four
 span all \(4\times4\) matrices.  Derive the duality formulas relating the
 rank-two, rank-three, and rank-four products to \(\gamma_5\), tracking
 their phases against the formulas printed in Appendix~C.
 \item Derive the vanishing of odd traces, the two- and four-gamma traces,
 the corresponding \(\gamma_5\) traces, and a recursion for longer traces.
 Prove the contraction identity for an antisymmetric rank-\(r\) product
 and recover the ordinary-string contractions displayed in Appendix~C.
 \item Solve \((\slashed p-m)u=0\) and \((\slashed p+m)v=0\) explicitly.
 Establish the covariant normalization and the two completeness relations,
 including the relation between the two-component spin bases.
\end{enumerate}
State every phase and normalization convention used in the derivation.
\end{exercise}
